\section{Symmetric and exterior powers}
\label{section-symmetric-exterior}

\noindent
Let $(X, \mathcal{O}_X)$ be a ringed space.
Let $\mathcal{F}$ be an $\mathcal{O}_X$-module.
We define the {\it tensor algebra} of $\mathcal{F}$ to be
the sheaf of noncommutative $\mathcal{O}_X$-algebras
$$
\text{T}(\mathcal{F})
=
\text{T}_{\mathcal{O}_X}(\mathcal{F})
= \bigoplus\nolimits_{n \geq 0} \text{T}^n(\mathcal{F}).
$$
Here $\text{T}^0(\mathcal{F}) = \mathcal{O}_X$,
$\text{T}^1(\mathcal{F}) = \mathcal{F}$
and for $n \geq 2$ we have
$$
\text{T}^n(\mathcal{F}) =
\mathcal{F} \otimes_{\mathcal{O}_X} \ldots \otimes_{\mathcal{O}_X} \mathcal{F}
\ \ (n\text{ factors})
$$
We define $\wedge(\mathcal{F})$ to be the quotient of
$\text{T}(\mathcal{F})$ by the two sided ideal generated by
local sections $s \otimes s$ of $\text{T}^2(\mathcal{F})$ where
$s$ is a local section of $\mathcal{F}$. This is
called the {\it exterior algebra} of $\mathcal{F}$.
Similarly, we define $\text{Sym}(\mathcal{F})$ to be
the quotient of $\text{T}(\mathcal{F})$ by the two
sided ideal generated by local sections of the form
$s \otimes t - t \otimes s$ of $\text{T}^2(\mathcal{F})$.

\medskip\noindent
The algebras $\text{T}(\mathcal{F})$, $\wedge(\mathcal{F})$, and
$\text{Sym}(\mathcal{F})$ are graded $\mathcal{O}_X$-algebras
in the sense of Differential Graded Sheaves, Definition
\ref{sdga-definition-ga}.
Moreover $\text{Sym}(\mathcal{F})$
is commutative, and $\wedge(\mathcal{F})$ is graded commutative.

\begin{lemma}
\label{lemma-local-tensor-algebra}
In the situation described above.
The sheaf $\wedge^n\mathcal{F}$ is the sheafification of the
presheaf
$$
U \longmapsto \wedge^n_{\mathcal{O}_X(U)}(\mathcal{F}(U)).
$$
See Algebra, Section \ref{algebra-section-tensor-algebra}.
Similarly, the sheaf $\text{Sym}^n\mathcal{F}$ is the sheafification
of the presheaf
$$
U \longmapsto \text{Sym}^n_{\mathcal{O}_X(U)}(\mathcal{F}(U)).
$$
\end{lemma}

\begin{proof}
Omitted. It may be more efficient to define $\text{Sym}(\mathcal{F})$
and $\wedge(\mathcal{F})$ in this way instead of the method
given above.
\end{proof}

\begin{lemma}
\label{lemma-stalk-tensor-algebra}
In the situation described above. Let $x \in X$.
There are canonical isomorphisms of $\mathcal{O}_{X, x}$-modules
$\text{T}(\mathcal{F})_x = \text{T}(\mathcal{F}_x)$,
$\text{Sym}(\mathcal{F})_x = \text{Sym}(\mathcal{F}_x)$, and
$\wedge(\mathcal{F})_x = \wedge(\mathcal{F}_x)$.
\end{lemma}

\begin{proof}
Clear from Lemma \ref{lemma-local-tensor-algebra} above, and
Algebra, Lemma \ref{algebra-lemma-colimit-tensor-algebra}.
\end{proof}

\begin{lemma}
\label{lemma-pullback-tensor-algebra}
Let $f : (X, \mathcal{O}_X) \to (Y, \mathcal{O}_Y)$ be a morphism of
ringed spaces. Let $\mathcal{F}$ be a sheaf of $\mathcal{O}_Y$-modules.
Then $f^*\text{T}(\mathcal{F}) = \text{T}(f^*\mathcal{F})$,
and similarly for the exterior and symmetric algebras associated
to $\mathcal{F}$.
\end{lemma}

\begin{proof}
Omitted.
\end{proof}

\begin{lemma}
\label{lemma-presentation-sym-exterior}
Let $(X, \mathcal{O}_X)$ be a ringed space.
Let $\mathcal{F}_2 \to \mathcal{F}_1 \to \mathcal{F} \to 0$
be an exact sequence of sheaves of $\mathcal{O}_X$-modules.
For each $n \geq 1$ there is an exact sequence
$$
\mathcal{F}_2 \otimes_{\mathcal{O}_X} \text{Sym}^{n - 1}(\mathcal{F}_1)
\to
\text{Sym}^n(\mathcal{F}_1)
\to
\text{Sym}^n(\mathcal{F})
\to
0
$$
and similarly an exact sequence
$$
\mathcal{F}_2 \otimes_{\mathcal{O}_X} \wedge^{n - 1}(\mathcal{F}_1)
\to
\wedge^n(\mathcal{F}_1)
\to
\wedge^n(\mathcal{F})
\to
0
$$
\end{lemma}

\begin{proof}
See Algebra, Lemma \ref{algebra-lemma-presentation-sym-exterior}.
\end{proof}

\begin{lemma}
\label{lemma-tensor-algebra-permanence}
Let $(X, \mathcal{O}_X)$ be a ringed space.
Let $\mathcal{F}$ be a sheaf of $\mathcal{O}_X$-modules.
\begin{enumerate}
\item If $\mathcal{F}$ is locally generated by sections,
then so is each $\text{T}^n(\mathcal{F})$,
$\wedge^n(\mathcal{F})$, and $\text{Sym}^n(\mathcal{F})$.
\item If $\mathcal{F}$ is of finite type,
then so is each $\text{T}^n(\mathcal{F})$,
$\wedge^n(\mathcal{F})$, and $\text{Sym}^n(\mathcal{F})$.
\item If $\mathcal{F}$ is of finite presentation,
then so is each $\text{T}^n(\mathcal{F})$,
$\wedge^n(\mathcal{F})$, and $\text{Sym}^n(\mathcal{F})$.
\item If $\mathcal{F}$ is coherent,
then for $n > 0$ each $\text{T}^n(\mathcal{F})$,
$\wedge^n(\mathcal{F})$, and $\text{Sym}^n(\mathcal{F})$
is coherent.
\item If $\mathcal{F}$ is quasi-coherent,
then so is each $\text{T}^n(\mathcal{F})$,
$\wedge^n(\mathcal{F})$, and $\text{Sym}^n(\mathcal{F})$.
\item If $\mathcal{F}$ is locally free,
then so is each $\text{T}^n(\mathcal{F})$,
$\wedge^n(\mathcal{F})$, and $\text{Sym}^n(\mathcal{F})$.
\end{enumerate}
\end{lemma}

\begin{proof}
These statements for $\text{T}^n(\mathcal{F})$ follow
from Lemma \ref{lemma-tensor-product-permanence}.

\medskip\noindent
Statements (1) and (2) follow from the fact that
$\wedge^n(\mathcal{F})$ and $\text{Sym}^n(\mathcal{F})$
are quotients of $\text{T}^n(\mathcal{F})$.

\medskip\noindent
Statement (6) follows from
Algebra, Lemma \ref{algebra-lemma-free-tensor-algebra}.

\medskip\noindent
For (3) and (5) we will use
Lemma \ref{lemma-presentation-sym-exterior} above.
By locally choosing a presentation
$\mathcal{F}_2 \to \mathcal{F}_1 \to \mathcal{F} \to 0$
with $\mathcal{F}_i$ free, or finite free and applying the
lemma we see that $\text{Sym}^n(\mathcal{F})$, $\wedge^n(\mathcal{F})$
has a similar presentation; here we use (6) and
Lemma \ref{lemma-tensor-product-permanence}.

\medskip\noindent
To prove (4) we will use
Algebra, Lemma \ref{algebra-lemma-present-sym-wedge}.
We may localize on $X$ and assume that
$\mathcal{F}$ is generated by a finite set
$(s_i)_{i \in I}$ of global sections.
The lemma mentioned above
combined with Lemma \ref{lemma-local-tensor-algebra} above
implies that for $n \geq 2$
there exists an exact sequence
$$
\bigoplus\nolimits_{j \in J}
\text{T}^{n - 2}(\mathcal{F})
\to
\text{T}^n(\mathcal{F})
\to
\text{Sym}^n(\mathcal{F})
\to
0
$$
where the index set $J$ is finite. Now we know that
$\text{T}^{n - 2}(\mathcal{F})$ is finitely generated
and hence the image of the first arrow is a coherent
subsheaf of $\text{T}^n(\mathcal{F})$, see Lemma \ref{lemma-coherent-abelian}.
By that same lemma we conclude that $\text{Sym}^n(\mathcal{F})$ is
coherent.
\end{proof}

\begin{lemma}
\label{lemma-whole-tensor-algebra-permanence}
Let $(X, \mathcal{O}_X)$ be a ringed space.
Let $\mathcal{F}$ be a sheaf of $\mathcal{O}_X$-modules.
\begin{enumerate}
\item If $\mathcal{F}$ is quasi-coherent,
then so is each $\text{T}(\mathcal{F})$,
$\wedge(\mathcal{F})$, and $\text{Sym}(\mathcal{F})$.
\item If $\mathcal{F}$ is locally free,
then so is each $\text{T}(\mathcal{F})$,
$\wedge(\mathcal{F})$, and $\text{Sym}(\mathcal{F})$.
\end{enumerate}
\end{lemma}

\begin{proof}
It is not true that an infinite direct sum $\bigoplus \mathcal{G}_i$ of
locally free modules is locally free, or that an
infinite direct sum of quasi-coherent modules
is quasi-coherent. The problem is that given a
point $x \in X$ the open neighbourhoods $U_i$ of $x$ on which $\mathcal{G}_i$
becomes free (resp.\ has a suitable presentation) may have an intersection
which is not an open neighbourhood of $x$. However, in the
proof of Lemma \ref{lemma-tensor-algebra-permanence} we saw that
once a suitable open neighbourhood for $\mathcal{F}$ has been chosen,
then this open neighbourhood works for each of the sheaves
$\text{T}^n(\mathcal{F})$, $\wedge^n(\mathcal{F})$ and
$\text{Sym}^n(\mathcal{F})$.
The lemma follows.
\end{proof}







